\section{Introduction}
\label{sec:intro}
A symbol-based augmentative and alternative communication (AAC) board lets a person compose a message by selecting pictograms.
The labels alone do not determine a sentence: \emph{Water} can request a drink, whereas \emph{No}\,+\,\emph{Water} can refuse one.
Adding grammatical structure must therefore preserve both the selected concepts and the communicative act.
This matters for users who rely on the device to express needs, discomfort, and refusals to a listener who may not see the board.

\begin{figure}[t]
\centering
\begin{tikzpicture}[font=\small,>=Latex,
 box/.style={draw=exampleblue!55,rounded corners=2pt,align=left,text width=6.65cm,inner sep=6pt},
 node distance=5mm]
\node[box,fill=panelgray] (a) {\textbf{Selection and message}\quad Water + request\\Concept: water; polarity: positive; speaker: user};
\node[box,below=of a] (b) {\textbf{A fluent candidate can add content}\\I want water \bad{tomorrow}.\quad\emph{[added: time]}};
\node[box,below=of b,fill=exampleblue!5] (c) {\textbf{Check, then choose a route}\\Model sentence $\rightarrow$ gate $\rightarrow$ checked template\\\good{I want water.}\quad or clarification / literal preview};
\draw[->,thick] (a)--(b);
\draw[->,thick] (b)--(c);
\end{tikzpicture}

\caption{Concept-gated realization, illustrated with an authored example. A water request must not acquire an unselected time or lose its act. Red text is labelled as added content; blue text is the intended realization. Acceptance is a heuristic check, not a semantic guarantee.}
\label{fig:teaser}
\end{figure}

For a water request, a single symbol can be realized as a complete sentence in each language.
The following are authored illustrations; recorded model outputs are kept separate in Appendix~\ref{app:traces}.
\begin{quote}\small
\lang{English reference}\quad\good{I want water.}\\[2pt]
\lang{Hindi}\quad\good{\hi{मुझे पानी चाहिए।}}\\[2pt]
\lang{Gujarati}\quad\good{\gu{મને પાણી જોઈએ છે.}}\\[2pt]
\lang{Marathi}\quad\good{\hi{मला पाणी हवे आहे.}}
\end{quote}
The three Indic sentences share the English gloss \emph{``I want water.''}
Feelings, refusals, and actions require other constructions, so a realization policy must distinguish them rather than place every label in a request frame.

A language model can supply these constructions but can also replace a concept or add an explanation (Figure~\ref{fig:teaser}).
Our approach builds a \emph{concept ledger} from the selection and checks each proposed sentence for concepts, polarity, act cues, and unrequested content before returning it.
The revised workflow also checks its template fallback and asks for a message type when the selection is ambiguous.
These checks are deliberately inspectable, but their finite vocabulary does not determine every possible meaning.

We evaluate what candidate checking adds beyond templates.
On the archived grid, the policy improves its ungated cascade but trails a template baseline.
We trace errors across input interpretation, generation, checking, and scoring.
Symbol-board users compose messages; caregivers, clinicians, and researchers configure vocabulary.
We make three contributions:
\begin{itemize}
\item \textbf{A concept-gated realization policy} with explicit concept, polarity, and act checks, and a revised checked-template fallback.
\item \textbf{A four-language interaction} that separates generated sentences, template sentences, clarification, and literal previews.
\item \textbf{A paired evaluation and error audit} that exposes the template advantage, interface failures, and scorer limitations, followed by a separately labelled repair regression.
\end{itemize}
The contribution is the design and evaluation of these checks for Indic symbol realization, rather than multilingual AAC or grammatical expansion alone.
